
\begin{exercise}
The value of a car, $V$, in dollars, is a function of the number of years, $t$, since it was purchased; that is, $V=g(t)$.
\begin{parts}
\item What is the independent variable?

\item In what units is the independent variable measured?

\item What is the dependent variable?

\item In what units is the dependent variable measured?
\end{parts}
\end{exercise}


\begin{solution}

\begin{parts}
\item $t$
\item years
\item $V$
\item dollars
\end{parts}
\end{solution}





\begin{exercise}
The amount of caffeine, $C$, measured in milligrams, in a person's body $t$ hours after drinking a cup of coffee is given by the function $C=f(t)$. What do each of the following statements tell you in practical terms; that is, in terms of time and caffeine? Answer in complete sentences and include units with each number.


\begin{parts}
\item $f(0)=96$
\item $f(5)=48$
\item $f(24)\approx 0$ 
\end{parts}
\end{exercise}

\begin{solution}

\begin{parts}
\item The amount of caffeine in a person's body $0$ hours after drinking a cup of coffee is 96 mg.
\item The amount of caffeine in a person's body $5$ hours after drinking a cup of coffee is 48 mg.
\item The amount of caffeine in a person's body $24$ hours after drinking a cup of coffee is approximately 0 mg.
\end{parts}
\end{solution}











\begin{exercise}
The total number of units, $P$, produced after $t$ hours of production is given by the function $P=h(t)$. What do each of the following statements tell you in practical terms; that is, in terms of units of product produced and time? Answer in complete sentences and include units with each number.


\begin{parts}
        \item  $h(1)=10$
        \item  $0=h(0)$
        \item  $h(3) \approx 29 $
\end{parts}
\end{exercise}

\begin{solution}

\begin{parts}
\item A total of 10 units of the product are produced in 1 hour.
\item A total of 0 units of the product are produced in 0 hours.
\item A total of approximately 29 units of the product are produced in 3 hours.
\end{parts}
\end{solution}









\begin{exercise}
The total cost of a meal in a restaurant, $C$, in dollars, as a function of the menu price of the meal, $p$, also in dollars, is given by
\[C=p+0.20p\]
where the term $0.20p$ corresponds to a 20\% tip.


\begin{parts}
\item What is the input variable and what units is it measured in? 

\item What is the output variable and what units is it measured in? 

\item Calculate the total cost of a meal whose menu price is $\$25$.
\end{parts}

\end{exercise}





\begin{solution}

\begin{parts}
\item $p$; dollars
\item $C$; dollars
\item $\$30$
\end{parts}
\end{solution}








\begin{exercise}
A bakery makes 6 batches of cookies an hour, with a batch typically consisting of 12 cookies. The total number of batches baked, $B$, is measured as function of time $t$, in hours, so that $B=f(t)=6t$.


\begin{parts}
\item What is the independent variable?
\item What is the dependent variable?
\item Evaluate $f(3)$. What is the practical meaning of $f(3)$ in the applied context of this problem?
\end{parts}

\end{exercise}





\begin{solution}

\begin{parts}
\item The independent variable is $t$, which represents the time spent baking in hours.
\item The dependent variable is $B$, which represents the number of batches of cookies made by the bakery after $t$ hours have passed.
\item $f(3)=18$; The bakery makes 18 batches of cookies (or approximately 216 cookies) in 3 hours.
\end{parts}
\end{solution}











\begin{exercise}
A driver is heading to a faraway town. The amount of fuel, $G$, in gallons, left in the fuel tank is a function of the number of miles $m$ driven during the trip; that is, $G=f(m)$.


\begin{parts}

\item What does the statement $f(70)=6$ tell you in practical terms? What are units of the numbers 70 and 6?

\item What does the statement  $f(200)=1$ tell you in practical terms?
\end{parts}
\end{exercise}



\begin{solution}

\begin{parts}
\item The amount of fuel left in the fuel tank 70 miles into the drive is 6 gallons.
\item The amount of fuel left in the fuel tank 200 miles into the drive is 1 gallon.
\end{parts}
\end{solution}








\begin{exercise}
Let $f(x)=2x+1$. Evaluate: 


\begin{parts}
\item $f(3)$
\item $f(-4.5)$
\item $f(0)$ 
\end{parts}
\end{exercise}



\begin{solution}

\begin{parts}
\item $7$ 
\item $-8$
\item $1$
\end{parts}
\end{solution}








\begin{exercise}
Let $f(x)=4-3x$. Evaluate and simplify if possible: 


\begin{parts}
\item $f(3b)$
\item $f(b+1)$
\item $f(b+3)-f(3)$
\end{parts}
\end{exercise}



\begin{solution}

\begin{parts}
\item $-9b+4$ 
\item $-3b+1$
\item $-3b$
\end{parts}
\end{solution}






\begin{exercise}
Let $g(x)=3x-x^2$. Evaluate: 


\begin{parts}
\item $g(2)$
\item $g(-2)$
\item $g(\sqrt{2})$
\end{parts}
\end{exercise}




\begin{solution}

\begin{parts}
\item $2$ 
\item $-10$
\item $3\sqrt{2}-2$
\end{parts}
\end{solution}







\begin{exercise}
Let $g(x)=x^2+3$. Evaluate and simplify if possible: 


\begin{parts}
\item $g(h-1)$
\item $g(h+1)$
\item $g(b+1)-g(1)$
\end{parts}
\end{exercise}


\begin{solution}

\begin{parts}
\item $h^2-2h+4$ 
\item $h^2+2h+4$
\item $b^2+2b$
\end{parts}
\end{solution}








\begin{exercise}
Let $\displaystyle f(x)=\frac{1}{x+1}$. Evaluate. If the value is undefined, say so. 


\begin{parts}
\item $f(0)$
\item $f(-1)$
\item $f(-2)$
\end{parts}
\end{exercise}


\begin{solution}

\begin{parts}
\item $1$ 
\item undefined
\item $-1$
\end{parts}
\end{solution}








\begin{exercise}
Let $\displaystyle f(x)=\frac{3}{x+2}$. Evaluate and simplify if possible. 


\begin{parts}
\item $f(h-1)$
\item $f(h+1)$
\item $f(b)-4f(1)$
\end{parts}
\end{exercise}


\begin{solution}

\begin{parts}
\item $\frac{3}{h+1}$ 
\item $\frac{3}{h+3}$
\item $\frac{3}{b+2}-4$
\end{parts}
\end{solution}





\begin{exercise}
Let $f(t)=t^2+\frac{1}{t}$. Evaluate and simplify if possible.


\begin{parts}
        \item  $f(2)$
        \item  $f(h+1)$
        \item  $f(h+1) - f(2)$
\end{parts}
\end{exercise}


\begin{solution}

\begin{parts}
\item $\frac{9}{2}$
\item $h^2+2h+1+\frac{1}{h+1}$
\item $h^2+2h+\frac{1}{h+1}-\frac{7}{2}$
\end{parts}
\end{solution}






\begin{exercise}
Let $g(x)=2x^2+\frac{3}{2}x$ and $f(x)=x+3$. Evaluate and simplify if possible.


\begin{parts}
\item  $f(2)+g(1)$
\item  $g(1)+2f(x+1)$
\item  $f(2x)-g(4+x)$
\end{parts}
\end{exercise}


\begin{solution}

\begin{parts}
\item $\frac{17}{2}$
\item $2x+\frac{23}{2}$
\item $-2x^2-\frac{31}{2}x-35$
\end{parts}
\end{solution}






\begin{exercise}
A URI student tutors for MTH 103. The cost of each session in terms of hours, $t$, is given by the function $c(t)=10t+15$.


\begin{parts}
\item What is the cost of a 3 hour session?
\item If $c(t)=\$60$, how many hours was the session?
\end{parts}
\end{exercise}


\begin{solution}

\begin{parts}
\item $\$45$
\item $4.5$ hours
\end{parts}
\end{solution}









\begin{exercise}
The value of a car, $V$, in dollars, $t$ years after purchase is given by the function $V=g(t)$, where $g(t)=16500-1500t$.


\begin{parts}
\item What is the value of the car at the time of purchase?

\item What is the value of the car $5$ years after purchase?

\item After how many years is the car worth nothing?
\end{parts}
\end{exercise}


\begin{solution}

\begin{parts}
\item $\$16500$
\item $\$9000$
\item $11$ years
\end{parts}
\end{solution}







\begin{exercise}
Since you see lightning immediately and it takes the sound of thunder about 5 seconds to travel a mile, you can calculate the distance between you and the lightning\footnote{https://www.weather.gov/safety/lightning-science-thunder, accessed: 6/26/20}. Count the number of seconds, $S$, between the flash of lightning and the sound of thunder. Then the distance, $D$, in miles, between you and the lightning is given by the function:
\[D=\frac{S}{5}\]


\begin{parts}
\item Identify the independent and the dependent variable.

\item How far is the lightning if you counted 4 seconds between the flash and the thunder?
\end{parts}
\end{exercise}

\begin{solution}

\begin{parts}
\item $S$ is the independent variable; $D$ is the dependent variable
\item $\frac{4}{5}$ of a mile
\end{parts}
\end{solution}







\begin{exercise}
The cost of a medium pizza is $\$10$ plus $\$1.75$ per topping. Write a function $f(x)$ to describe the cost of a medium pizza with $x$ number of toppings.
\end{exercise}

\begin{solution}
$f(x)=1.75x+10$
\end{solution}







\begin{exercise}
Find a formula for a function $h(x)$ for each of the following scenarios:


\begin{parts}
\item $h$ multiplies the input by 7 then adds $5$ to the result.

\item $h$ adds $5$ to the input and then multiplies the result by 7.
\end{parts}

\noindent Are the functions in (a) and in (b) the same?
\end{exercise}


\begin{solution}

\begin{parts}
\item $7x+5$
\item $7(x+5)$
\end{parts}
\noindent The functions are not the same.
\end{solution}






\begin{exercise}
Find a formula for a function $f(x)$ for each of the following scenarios:


\begin{parts}
\item $f$ takes the square of the input, multiplies the result by 5, then subtracts 8.

\item $f$ multiplies the input by 5, takes the square of the result, then subtracts 8.

\item  $f$ subtracts 8 from the input, takes the square of the result, then multiplies by 5.
\end{parts}

\noindent Are any of the three functions the same?
\end{exercise}


\begin{solution}

\begin{parts}
\item $5x^2-8$
\item $(5x)^2-8 = 25x^2-8$
\item $5(x-8)^2$
\end{parts}
\noindent None of the three functions are the same.
\end{solution}










\begin{exercise}
Find a formula for a function $g(x)$ for each of the following scenarios:


\begin{parts}
\item $g$ takes the input, divides it by 3, then adds this to the square root of the input times 4.
\item $g$ takes the reciprocal of the square root of the input, adds $1$ to this, and then adds the input squared.
\end{parts}
\end{exercise}


\begin{solution}

\begin{parts}
\item $\frac{x}{3} + \sqrt{x}\cdot 4$
\item $\frac{1}{\sqrt{x}}+1+x^2$
\end{parts}
\end{solution}





\begin{exercise}
Find the domain of each of the following functions.


\begin{parts}[itemsep=3pt]
\item $\displaystyle f(x)=x-1$
\item $\displaystyle f(x)=\frac{1}{x-1}$
\item $\displaystyle f(x)=\frac{2}{x^2-4}$
\item $\displaystyle f(x)=\sqrt{x}$
\item $\displaystyle f(x)=x^2-4$  
\end{parts}
\end{exercise}


\begin{solution}

\begin{parts}
\item all real numbers or $(-\infty,\infty)$ in interval notation
\item all real numbers $x\neq1$ or $(-\infty,1)\cup(1,\infty)$ in interval notation
\item all real numbers $x\neq-2$, $x\neq2$ or $(-\infty,-2)\cup(-2,2)\cup(2,\infty)$ in interval notation
\item all real numbers $x\geq0$ or $[0,\infty)$ in interval notation
\item all real numbers or $(-\infty,\infty)$ in interval notation
\end{parts}
\end{solution}







\begin{exercise}
Find the domain of each of the following functions.


\begin{parts}
\item $f(x)=x^n$, where $n$ is a positive integer
\item $f(x)=x^2+x+\frac{2}{x}$
\item $f(x)=\frac{1}{\sqrt{x}}$
\end{parts}
\end{exercise}


\begin{solution}

\begin{parts}
\item all real numbers or $(-\infty,\infty)$ in interval notation
\item all real numbers $x\neq0$ or $(-\infty,0)\cup(0,\infty)$ in interval notation
\item all real numbers $x> 0$ or $(0,\infty)$ in interval notation
\end{parts}
\end{solution}










